\begin{papercard}{spencer2000}
\cf{Stages} \stagetag{W}. A practitioner modification of the Wharton et al.\ cognitive walkthrough for large development teams: the \emph{streamlined CW}, with two questions instead of four, ground rules and lighter recording, plus a rough external validity check.

\cf{Problem} The Wharton CW failed in a large software company because of three social constraints: time pressure, lengthy design discussions and design defensiveness \pc{spencer2000}{2}{Social Constraints That Hampered the Effectiveness of the}. Recording a plausible story for every step produced voluminous, low-value output, and asking four questions per obviously fine step felt slow \pc{spencer2000}{2}{Voluminous output of WCW sessions}. Teams struggled to distinguish questions 1--3.

\cf{Modified process} Define inputs as in Wharton. Convene: describe goals and what will \emph{not} be done, explicitly defuse defensiveness, post ground rules, assign roles, ask for submission to the leader's facilitation \pc{spencer2000}{3}{Explicitly defuse defensiveness}. Present findings in softer language (``a potential problem'') to match the method's tentative nature \pc{spencer2000}{3}{take care to use softer}. Train with two worked examples (one passing, one failing) and model the data format.

\cf{Two questions} (1) Will the user know what to do at this step? (2) If the user does the right thing, will they know they did, and are they making progress? \pc{spencer2000}{3}{2 questions from the streamlined CW}. Wharton's questions 1--3 are collapsed into the first \pc{spencer2000}{4}{questions 1-3 evaluate whether the user}.

\cf{Ground rules} No designing; no defending a design; no debating cognitive theory; the usability specialist leads \pc{spencer2000}{4}{Ground rules for conducting a streamlined CW}.

\cf{Recording} If both questions get a plausible story, nothing is recorded. Otherwise record the failure plus the missing user knowledge. Four data types: potential learnability problems, design ideas (marked ``DI''), design gaps, task-analysis flaws, with roles assigned for each \pc{spencer2000}{4}{If, for a particular action sequence, a plausible story is told}. Rephrase each bullet as a hypothesis to obtain consensus \pc{spencer2000}{5}{When a bullet is being captured, it is helpful for the}. Do not retool a faulty task analysis during the session \pc{spencer2000}{5}{important to not try to retool the task analysis during the}.

\cf{Trade-offs acknowledged} Collapsing questions gives coarser grain and less causal explanation. Banning design talk trades solutions for coverage. Not recording stories loses design rationale \pc{spencer2000}{5}{This modification to the CW method probably leads to a}.

\cf{Evaluation} One session of about 1.5 hours covered 32 action sequences with 8 people (3 usability, 1 designer, 4 PMs), after a 25-hour task analysis. Results: 24 potential problems (14 knowledge, 10 feedback) and 11 design ideas \pc{spencer2000}{6}{Twenty-four potential problems and l I design ideas were}. A later usability test found 9 problems, 6 of which the CW had identified. Of 13 CW predictions, 7 were directly confirmed, 4 indirectly and 2 not related (false positives) \pc{spencer2000}{6}{For the portions of the user interface covered by both the}.

\cf{Use for this thesis} A precedent for \emph{documented simplification} of CW. If the thesis uses fewer than four questions, cite Spencer, state the trade-off (coarser causes), and keep recording the missing user knowledge. Its separation of problems, design ideas, gaps and task-analysis flaws is a clean recording schema. ``Note and move on'' prevents analysis drift, analogous to staying on the correct path in Wharton. The softer-language guidance matches reporting harm-communication findings as \emph{potential} barriers.
\end{papercard}
